\section{Posterior stability and lifted discretization}
\label{sec:wellposed}

\subsection{Existence and stability in the data}

The first theorem is a property of Bayes' formula and does not require
the Girsanov representation.

\begin{theorem}[Posterior well-posedness]
\label{thm:wellposed}
Let $\G:\HH\to\R^k$ be bounded and linear, let $\Gamobs$ be positive
definite, and fix $\sigobs>0$.  If
\begin{equation}
  \int_{\HH}\norm{u}_{\HH}^2\,\muz(\mathrm du)<\infty ,
  \label{eq:prior-second-moment}
\end{equation}
then $0<Z(y)\leq1$ for every $y\in\R^k$, so
\eqref{eq:posterior} defines a probability measure.  Moreover, for
every $R<\infty$ there is a finite constant $C_R$ such that
\begin{equation}
  \dH(\mu^y,\mu^{y'})
  \leq C_R\norm{y-y'}_{\R^k},
  \qquad \norm{y},\norm{y'}\leq R .
  \label{eq:data-stability}
\end{equation}
The constant may depend on
$R,\sigobs,\G,\Gamobs$, and $\muz$.
\end{theorem}

\Cref{ass:drift} implies \eqref{eq:prior-second-moment} by
\Cref{lem:sde-moments}.  The role of \Cref{ass:girsanov} is different:
it gives the density \eqref{eq:terminal-density}, which is useful when
transferring local mass estimates from the Gaussian reference prior.
Neither Gaussianity nor absolute continuity is needed for
\Cref{thm:wellposed}.

\subsection{A full-space spectral lift}

Define
\begin{equation}
  \driftN(t,u)=P_N\drift(t,P_Nu),\qquad
  a_N(t,u)=\mathcal C^{-1/2}\driftN(t,u).
  \label{eq:projected-drift}
\end{equation}
Because $P_N$ commutes with $\mathcal C$, the projected drift remains
Cameron--Martin-valued and satisfies
\Cref{ass:drift} with constants independent of $N$.  Let $U^N$ solve
\begin{equation}
  \mathrm dU_t^N
  =\bigl[-\tfrac12U_t^N+\driftN(t,U_t^N)\bigr]\,\mathrm dt
   +\mathcal C^{1/2}\,\mathrm dW_t,\qquad
  U_0^N\sim\mu_{\mathrm{ref}},
  \label{eq:lifted-prior-sde}
\end{equation}
and let $\mu_{0,N}=\mathcal L(U_T^N)$.  Notice that the noise in
\eqref{eq:lifted-prior-sde} is not projected.  The high-frequency
component $(I-P_N)U_T^N$ therefore retains its reference Gaussian law.

The likelihood is approximated by
\begin{equation}
  \Phim_{\sigobs,N}(u;y)
  =\Phim_{\sigobs}(P_Nu;y),
  \label{eq:projected-potential}
\end{equation}
and the lifted posterior is
\begin{equation}
  \frac{\mathrm d\muyN}{\mathrm d\mu_{0,N}}(u)
  =\frac{\exp\{-\Phim_{\sigobs,N}(u;y)\}}{Z_N(y)} .
  \label{eq:lifted-posterior}
\end{equation}
Both $\muyN$ and $\muy$ are probability measures on the same space
$\HH$.

Let $\mathbb Q^\theta$ and $\mathbb Q_N^\theta$ denote the path
measures obtained from the reference OU law by the stochastic
exponentials associated with $a$ and $a_N$, respectively.

\begin{assumption}[Drift approximation]
\label{ass:drift-approximation}
For every $N$, the two stochastic exponentials are true martingales,
$\mathbb Q^\theta\ll\mathbb Q_N^\theta$, and the Girsanov entropy
identity holds:
\begin{equation}
  \KL(\mathbb Q^\theta\,\|\,\mathbb Q_N^\theta)
  =\frac12\,
    \mathbb E_{\mathbb Q^\theta}
      \int_0^T\norm{a(t,X_t)-a_N(t,X_t)}_{\HH}^2\,\mathrm dt
  =\frac12\varepsilon_N<\infty .
  \label{eq:drift-energy}
\end{equation}
\end{assumption}

This condition is a statement about the actual path measures.  It is
not inferred from a training loss unless a separate argument connects
that loss to the expectation in \eqref{eq:drift-energy}.

\begin{theorem}[Lifted discretization bound]
\label{thm:discrete}
Suppose \Cref{ass:drift,ass:girsanov,ass:drift-approximation} hold.
Let $\G$ be bounded and linear, let $\Gamobs$ be positive definite, and
fix $\sigobs>0$.  Then, for every $R<\infty$, there is a constant
$C_R<\infty$, independent of $N$, such that
\begin{equation}
  \sup_{\norm{y}\leq R}
  \dH(\muyN,\muy)
  \leq C_R\left(\sqrt{\varepsilon_N}+r_N\right),
  \qquad
  r_N^2=\sum_{j>N}\lambda_j .
  \label{eq:discretization-bound}
\end{equation}
Consequently, $\dH(\muyN,\muy)\to0$ uniformly on bounded data sets
whenever $\varepsilon_N\to0$.
\end{theorem}

The two terms have distinct origins.  Data processing for relative
entropy controls the terminal prior error by the pathwise drift energy.
The likelihood error is controlled by the unresolved Gaussian
component, whose mean-square norm is exactly $r_N^2$.  The proof is in
\Cref{sec:appendix-discretization}.

\begin{corollary}[Polynomial and geometric covariance tails]
\label{cor:tail-rates}
Under the assumptions of \Cref{thm:discrete}, suppose
$\varepsilon_N=O(N^{-2q})$ for some $q>0$.
\begin{enumerate}
  \item If $\lambda_j\leq c j^{-2s}$ for some $s>1/2$, then
  \[
    \dH(\muyN,\muy)
    =O\!\left(N^{-\min\{q,s-1/2\}}\right)
  \]
  uniformly for $\norm{y}\leq R$.
  \item If $\lambda_j\leq c e^{-\alpha j}$ for some $\alpha>0$, then
  \[
    \dH(\muyN,\muy)
    =O\!\left(N^{-q}+e^{-\alpha N/2}\right)
  \]
  uniformly for $\norm{y}\leq R$.
\end{enumerate}
\end{corollary}

\begin{proof}
For polynomial decay, the integral test gives
$\sum_{j>N}\lambda_j=O(N^{1-2s})$.  For geometric decay, summing the
geometric series gives
$\sum_{j>N}\lambda_j=O(e^{-\alpha N})$.  Substitute these estimates
and $\sqrt{\varepsilon_N}=O(N^{-q})$ into
\eqref{eq:discretization-bound}.
\end{proof}

\begin{proposition}[Why the lift is necessary]
\label{prop:singular-collapse}
Suppose $\nu_N$ is supported on $\HN$ and $\muy\ll\muz\ll
\mu_{\mathrm{ref}}$.  Then $\nu_N\perp\muy$ and
\begin{equation}
  \dH(\nu_N,\muy)=1 .
\end{equation}
\end{proposition}

\begin{proof}
Every proper finite-dimensional subspace of $\HH$ has
$\mu_{\mathrm{ref}}$-measure zero.  Hence
$\muy(\HN)=0$, whereas $\nu_N(\HN)=1$.  The measures are mutually
singular, and \eqref{eq:hellinger} gives Hellinger distance one.
\end{proof}

Thus \Cref{thm:discrete} does not assert Hellinger convergence of a
collapsed finite-dimensional posterior.  Such approximations may
converge after projection, or in a weaker metric, but that is a
different statement.
